\Needspace{12\baselineskip}
\section{Positive volume without bigness}
\label{sec:radial-positive-volume-nonbig}

A nef boundary can retain positive fixed-density volume while losing every strict current representative. A radial potential on a projective base compensates the vanishing coefficient in a fiber direction. The resulting determinant controls all Gauduchon quotients, although a closed semipositive test still obstructs bigness.

Fix an integer $m\geq2$, put $n=m+1$ and $t=(m-1)/m$, and let $B=\mathbb P^m$. Give $B$ and $\mathbb P^1$ their Fubini--Study forms $\omega,w$, normalized by $\int_B\omega^m=\int_{\mathbb P^1}w=1$. On $B$ define
\[
 x([Z_0:\cdots:Z_m])=
 \frac{\sum_{j=1}^m|Z_j|^2}{\sum_{j=0}^m|Z_j|^2},
 \qquad a=x/t,\qquad p=[1:0:\cdots:0].
\]
On $X=B\times\mathbb P^1$ use pullbacks implicitly and set
\[
 h=\omega+w,\qquad \mu=\omega^m\wedge w,
 \qquad\theta=\omega+a w.
\]
For $\epsilon>0$ put
\begin{equation}\label{eq:radial-volume-normalization}
 H_\epsilon=m\int_0^1\frac{y^{m-1}}{y/t+\epsilon}\,\mathrm{d}y.
\end{equation}

\begin{lemma}\label{lem:radial-volume-representatives}
There are smooth real $\varphi_\epsilon$ on $B$ such that $A_\epsilon=\omega+\ddc\varphi_\epsilon>0$ and
\begin{equation}\label{eq:radial-volume-smooth-density}
 A_\epsilon^m=\frac{\omega^m}{H_\epsilon(a+\epsilon)}.
\end{equation}
They converge uniformly to a bounded continuous $\varphi\in\PSH(B,\omega)$, smooth off $p$. On $B\setminus\{p\}$ its current $A=\omega+\ddc\varphi$ is smooth strictly positive and satisfies
\begin{equation}\label{eq:radial-volume-endpoint-density}
 A^m=a^{-1}\omega^m.
\end{equation}
Moreover $H_\epsilon\to1$ as $\epsilon\downarrow0$.
\end{lemma}
\begin{proof}
In the affine chart $Z_0=1$, write $s=|z|^2$. Then $x=s/(1+s)$ and $\omega=\ddc\log(1+s)$. Polar integration gives
\begin{equation}\label{eq:radial-volume-integration}
 \int_B b(x)\omega^m=m\int_0^1 b(y)y^{m-1}\,\mathrm{d}y
\end{equation}
for every integrable radial function: the normalized density in $s$ is $m s^{m-1}(1+s)^{-m-1}\,\mathrm{d}s$. In particular $\int_B a^{-1}\omega^m=mt/(m-1)=1$. Monotone convergence in \eqref{eq:radial-volume-normalization} gives $H_\epsilon\to1$.

Use the positive root to define
\[
 q_\epsilon(x)^m=\frac m{H_\epsilon}
       \int_0^x\frac{y^{m-1}}{y/t+\epsilon}\,\mathrm{d}y,
 \qquad
 \varphi_\epsilon(x)=\int_0^x
       \frac{q_\epsilon(y)-y}{y(1-y)}\,\mathrm{d}y.
\]
At zero,
\[
 q_\epsilon(x)=x\left(\frac m{H_\epsilon}
       \int_0^1\frac{u^{m-1}}{xu/t+\epsilon}\,\mathrm{d}u\right)^{1/m},
\]
so it is $x$ times a smooth positive function. At one $q_\epsilon(1)=1$. Both apparent singularities of $\varphi_\epsilon'$ are removable smoothly. Composing with the global smooth projective function $x$ gives a smooth potential on $B$.

For a radial local potential $F(s)$, setting $q=sF'(s)$ gives tangential Hessian eigenvalue $q/s$ and radial eigenvalue $\mathrm{d}q/\mathrm{d}s$. For $F_\epsilon=\log(1+s)+\varphi_\epsilon(x)$ one has $sF_\epsilon'=q_\epsilon(x)$. Differentiating its defining integral shows both eigenvalues positive, including at $s=0$ by the smooth factorization. Their product relative to the Fubini--Study determinant is
\[
 \frac{A_\epsilon^m}{\omega^m}
 =\frac{q_\epsilon^{m-1}q_\epsilon'}{x^{m-1}}
 =\frac1{H_\epsilon(a+\epsilon)}.
\]
At every point of the infinity hyperplane, the global smooth form is semipositive as a limit of its positive affine matrices. The displayed determinant extends with a strictly positive value there, so every eigenvalue is positive. This proves global strict positivity and \eqref{eq:radial-volume-smooth-density}.

For $x>0$ the flux tends to $x^t$. Its normalized cumulative density is a decreasing reweighting of $m y^{m-1}\,\mathrm{d}y$, and an increasing reweighting of $(m-1)y^{m-2}\,\mathrm{d}y$: the latter weight ratio is proportional to $y/(y+t\epsilon)$. Comparing the averages below and above $x$ gives
\[
 x^m\leq q_\epsilon(x)^m\leq x^{m-1}.
\]
Consequently $x\leq q_\epsilon(x)\leq x^t$. The function
\[
 \varphi(x)=\int_0^x\frac{y^t-y}{y(1-y)}\,\mathrm{d}y
\]
is bounded continuous, since $t>0$, and smooth off $p$; its derivative extends smoothly at one. The integrable function $(y^t-y)/[y(1-y)]$ bounds the derivatives of $\varphi_\epsilon$. Dominated convergence gives uniform convergence of the potentials on all $B$. Distributional convergence makes $A$ positive and closed. On each coordinate chart, a smooth local potential for $\omega$ plus $\varphi$ is a continuous uniform limit of the corresponding psh functions with $\varphi_\epsilon$, including the value at $p$. Thus it is genuinely $\omega$-psh. Off $p$ the limiting flux gives
\[
 \frac{A^m}{\omega^m}
 =\frac{(x^t)^{m-1}t x^{t-1}}{x^{m-1}}=\frac t x=\frac1a.
\]
Smoothness and the same positive-determinant argument include the infinity hyperplane. This proves \eqref{eq:radial-volume-endpoint-density}.
\end{proof}

\Needspace{14\baselineskip}
\begin{theorem}\label{thm:radial-positive-volume-nonbig}\index{Gauduchon variational volume!positive nonbig background}
The background $\theta$ on $X$ is smooth, semipositive, nef and nonclosed. It is not big, but
\begin{equation}\label{eq:radial-positive-volume-nonbig}
 \lambda_\mu(\theta)=n^{1/n},\qquad\mathcal V_\mu(\theta)=n.
\end{equation}
For $\theta_\epsilon=\omega+(a+\epsilon)w$,
\[
 \mathcal V_\mu(\theta_\epsilon)=\frac n{H_\epsilon}.
\]
The smooth positive metrics
\begin{equation}\label{eq:radial-volume-minimizing-metrics}
 \gamma_\epsilon=(a+\epsilon)^{-1/m}
                 \bigl(A_\epsilon+(a+\epsilon)w\bigr)
\end{equation}
are balanced and minimize the defining quotient at $\theta_\epsilon$.
\end{theorem}
\begin{proof}
The coefficient $a$ is nonnegative, and $\theta+\eta h>0$ for every $\eta>0$, proving semipositivity and nefness. Since $a$ is nonconstant, $\mathrm{d}\theta=\mathrm{d}a\wedge w\ne0$.

Suppose that $\theta+\ddc v\geq\delta h$ for a quasi-plurisubharmonic $v$ and $\delta>0$. Choose a nonzero smooth $f\geq0$ on $B$, supported where $a<\delta$. The pullback $S=f\omega^m$ is smooth semipositive and closed by base top degree. Positivity and distributional Stokes would give
\[
 0\leq\int_X(\theta+\ddc v-\delta h)\wedge S
   =\int_B f(a-\delta)\omega^m<0.
\]
This contradiction excludes bigness; any Hermitian form is uniformly comparable to $h$ on compact $X$.

Pull back the bounded potential in \cref{lem:radial-volume-representatives}. The positive current $\Omega=A+a w=\theta+\ddc\varphi$ is smooth strictly positive off $Z=\{p\}\times\mathbb P^1$, where
\[
 \Omega^n=n a A^m\wedge w=n\mu.
\]
For every smooth positive Gauduchon metric $\gamma$, distributional Stokes replaces $\theta$ by $\Omega$ in its numerator. Since $Z$ has zero smooth volume, the matrix arithmetic--geometric mean inequality off $Z$ and positivity of $\Omega\wedge\gamma^{n-1}$ on $Z$ give
\[
 \int_X\theta\wedge\gamma^{n-1}
 \geq n^{1/n}\int_X
            (\mu/\gamma^n)^{1/n}\gamma^n
 =n^{1/n}D_\mu(\gamma).
\]
This proves the lower bound for every test in \cref{def:gaud-variational-volume}; no current power at $Z$ is used.

Put $c_\epsilon=n/H_\epsilon$ and $\Omega_\epsilon=A_\epsilon+(a+\epsilon)w$. Then $\Omega_\epsilon^n=c_\epsilon\mu$. The metrics in \eqref{eq:radial-volume-minimizing-metrics} have
\[
 T_\epsilon=\gamma_\epsilon^m
 =(a+\epsilon)^{-1}A_\epsilon^m
       +m A_\epsilon^{m-1}\wedge w.
\]
The first term is a base top-degree form; the second is closed. Thus $T_\epsilon$ is closed, proving the balanced and Gauduchon properties. Smooth Stokes and \eqref{eq:radial-volume-integration} give
\[
 \int_X\theta_\epsilon\wedge T_\epsilon
 =c_\epsilon\int_X(a+\epsilon)^{-1}\mu=n,
 \qquad
 D_\mu(\gamma_\epsilon)=c_\epsilon^{m/n}H_\epsilon.
\]
Their quotient is $c_\epsilon^{1/n}$. Applying the same matrix inequality to the smooth representative $\Omega_\epsilon$ gives the matching lower bound for every Gauduchon test, proving minimization and the perturbed volume formula. Since $\theta\leq\theta_\epsilon$, these actual metrics also give $\lambda_\mu(\theta)\leq c_\epsilon^{1/n}$. Letting $\epsilon\downarrow0$ proves \eqref{eq:radial-positive-volume-nonbig}. Its value agrees with Hermitian regularization by \cref{thm:gaud-volume-nef-limit}.
\end{proof}

The exact value above refers to the displayed probability $\mu$. Positivity persists for every smooth positive probability by \cref{prop:gaud-volume-density-comparison}. Thus positive fixed-density volume does not characterize bigness for nef nonclosed backgrounds, even on a projective manifold. The construction requires $m\geq2$ for its bounded endpoint potential. It leaves the surface case and the closed nef problem on general compact complex manifolds separate.
